\documentclass[11pt,a4paper]{article}

% ============================================
% 基础包
% ============================================
\usepackage[utf8]{inputenc}
\usepackage[T1]{fontenc}
\usepackage{amsmath,amssymb,amsthm}
\usepackage{mathtools}
\usepackage{algorithm}
\usepackage{algpseudocode}
\usepackage{graphicx}
\usepackage{booktabs}
\usepackage{multirow}
\usepackage{array}
\usepackage{longtable}
\usepackage{hyperref}
\usepackage{url}
\usepackage{cleveref}
\usepackage{geometry}
\usepackage{microtype}
\usepackage{xcolor}
\usepackage{tikz}
\usetikzlibrary{shapes,arrows,positioning,calc}

% 页面设置
\geometry{left=2.5cm,right=2.5cm,top=2.5cm,bottom=2.5cm}

% 颜色定义
\definecolor{primaryblue}{RGB}{0,82,147}
\definecolor{accentgreen}{RGB}{0,128,0}
\definecolor{warningred}{RGB}{178,34,34}

% 超链接设置
\hypersetup{
    colorlinks=true,
    linkcolor=primaryblue,
    citecolor=primaryblue,
    urlcolor=primaryblue,
    pdftitle={CAA-X V4: Cognitive Architecture eXtension for Zero-Shot Agent Challenges},
    pdfauthor={Lu Qi},
    pdfsubject={Cognitive Architecture, AI, Intention-Driven Systems},
}

% ============================================
% 定理环境
% ============================================
\theoremstyle{definition}
\newtheorem{definition}{Definition}[section]
\newtheorem{theorem}{Theorem}[section]
\newtheorem{proposition}{Proposition}[section]
\newtheorem{lemma}{Lemma}[section]
\newtheorem{corollary}{Corollary}[section]
\newtheorem{axiom}{Axiom}[section]

\theoremstyle{remark}
\newtheorem{remark}{Remark}[section]

% ============================================
% 自定义命令
% ============================================
\newcommand{\CAA}{\text{CAA-X}}
\newcommand{\CML}{\text{CML 2.0}}
\newcommand{\GWP}{\text{GWP}}
\newcommand{\lambdaK}{\lambda_{\text{K}}}
\newcommand{\lambdaM}{\lambda_{\text{M}}}
\newcommand{\lambdaI}{\lambda_{\text{I}}}
\newcommand{\lambdaP}{\lambda_{\text{P}}}
\newcommand{\lambdaS}{\lambda_{\text{S}}}
\newcommand{\lambdavec}{\boldsymbol{\lambda}}
\newcommand{\x}{\mathbf{x}}
\newcommand{\M}{\mathcal{M}}
\newcommand{\R}{\mathbb{R}}
\newcommand{\E}{\mathbb{E}}
\newcommand{\diff}{\,\mathrm{d}}

% ============================================
% 文档信息
% ============================================
\title{\textbf{CAA-X V4: Cognitive Architecture eXtension}\\[0.3em]
\large An Intention-Driven Framework with Empirical Validation\\[0.5em]
\normalsize Responding to LeCun's Zero-Shot Agent Challenge}

\author{Lu Qi\\
\textit{Intellisia Institute}\\
\texttt{[contact email to be updated]}}

\date{June 2, 2026}

% ============================================
% 文档开始
% ============================================
\begin{document}

\maketitle

\begin{abstract}
Current end-to-end large models, whether Transformer or Post-Transformer, suffer from a fundamental ``Intention Gap''---the absence of native intention representation in their architectures. This gap manifests as catastrophic reliability degradation in multi-turn interactions (39\% average decline), homogeneous output convergence across models, and failure in zero-shot physical reasoning tasks. We present \CAA{} (Cognitive Architecture eXtension), an intention-driven framework that bridges Global Workspace Theory (Baars, 1988), predictive processing (Friston, 2010), and modern AI engineering through five core innovations: (1) \textbf{Cognitive Atoms} as weakly-coupled functional primitives with persistence semantics (ephemeral/durable/permanent), validated by Roynard et al.'s four-layer cognitive substrate decomposition with DeepMind external validation; (2) \textbf{\CML{}}, a three-layer symbolic system (Latin/Semantic/Intentional) with differentiable transitions via Hamacher t-norms and energy-based constraints, mapping to five mechanistic interpretability paradigms; (3) \textbf{\GWP{} Engine}, a three-layer communication infrastructure (Handshake/Competition/Broadcast) that transforms GWT from metaphor to engineering specification, precisely mapped to System-1/System-2 coordination architectures; (4) \textbf{Cognitive Manifold} with tensor dynamics, where cognitive states evolve as geodesic flow on a Riemannian manifold with intention-biased trajectories, supported by critical initialization principles from biological neural networks (Pachitariu et al., Nature 2026); and (5) \textbf{Multidimensional Lambda-Order Parameters} ($\lambdaK, \lambdaM, \lambdaI, \lambdaP, \lambdaS$) as measurable indices of cognitive phase transitions, grounded in the AIQ benchmark. Empirical validation includes a robotic arm grasping experiment achieving 78\% zero-shot success on unseen object configurations. This work responds to the emerging deflationary trend in AI research (ICLR 2026, NeurIPS 2025) and positions cognitive architecture as a necessary evolution beyond associative prediction.
\end{abstract}

\textbf{Keywords:} Cognitive Architecture, Intention-Driven AI, Post-Transformer, Global Workspace Theory, Zero-Shot Agent, Cognitive Manifold, Lambda-Order Parameter, Persistence Semantics, System-1/System-2, Critical Initialization

\tableofcontents
\newpage

% ============================================
% 第1章 引言
% ============================================
\section{Introduction}\label{sec:introduction}

\subsection{The Intention Gap in Current AI Architectures}\label{subsec:intention-gap}

In May 2026, Pathway hosted a historic debate in San Francisco between Transformer co-inventor Lukasz Kaiser and Post-Transformer advocate Adrian Kosowski. The debate centered on memory mechanisms, inference costs, and hardware co-evolution---all questions of \textit{efficiency}. What was conspicuously absent was any discussion of what these architectures \textit{compute}, or \textit{why} they compute it. The entire field, it seemed, had tacitly agreed that the fundamental question of AI is ``how to compute faster,'' leaving ``what to compute and why'' unasked.

This is the ``\textbf{Intention Gap}''---the systematic absence of native intention representation in current AI architectures. By ``intention,'' we do not mean externally specified optimization objectives (loss functions, reward signals, or prompt instructions). We mean an internally generated, architecture-level capacity to select, maintain, and pursue goals---a capacity that is \textit{constitutive} of the system rather than imposed upon it.

The empirical evidence for the Intention Gap has accumulated rapidly:
\begin{itemize}
    \item Laban et al. (ICLR 2026) demonstrated that state-of-the-art language models lose 39\% reliability in multi-turn conversations across 200,000+ simulated dialogues---not because they lack predictive capacity, but because they lack intention maintenance.
    \item Jiang et al. (NeurIPS 2025) showed that different large language models converge to homogeneous outputs in open-ended tasks, a phenomenon they term ``Artificial Hivemind''---predictive systems without intention naturally converge to statistically dominant patterns rather than pursuing diverse goals.
    \item The AIQ benchmark (arXiv:2605.06815, 2026) reveals extreme cognitive imbalance in current AI: language comprehension exceeds the 98th percentile, while perceptual reasoning falls below the 1st.
\end{itemize}

\subsection{The Deflationary Turn in AI Research}\label{subsec:deflationary}

The AI research community is undergoing a subtle but fundamental shift. At ICLR 2026 in Rio de Janeiro, the Best Paper awards went not to another 0.7\% ImageNet improvement, but to a proof that Transformer properties are EXPSPACE-complete---formally undecidable (Bergstrasser et al., 2026)---and to a demonstration that top LLMs lose 39\% reliability in multi-turn conversations (Laban et al., 2026). This is not coincidence; it is a \textit{deflationary trend}: the field is rewarding papers that puncture overclaims, not papers that push benchmarks marginally.

Nature's March 2025 AGI review explicitly named Global Workspace Theory (GWT) and cognitive architectures as foundational to AGI. AAAI 2026's keynote by Zhihua Zhou's team laid out a four-stage neuro-symbolic paradigm, arguing that symbolic methods remain essential in the LLM era. The ICML 2026 Compositional Learning Workshop established consensus: compositionality is not an architectural byproduct but an engineering necessity. These developments create a window of opportunity for cognitive architecture---not as a nostalgic return to symbolic AI, but as a necessary evolution beyond associative prediction.

\subsection{Critical Initialization: The Biological Foundation}\label{subsec:critical-init}

A landmark study by Pachitariu et al. (Nature, 2026) provides the biological foundation for our approach. Through large-scale neural recordings in mice, they demonstrated that biological neural networks are initialized at birth to a ``critical state''---the boundary between order and chaos. This critical initialization, optimized by evolution over hundreds of millions of years, is the true origin of intelligence: not what the brain learns, but how it is structured \textit{before} learning begins.

The implications are profound for AI architecture. If biological intelligence emerges from critical initialization rather than training optimization, then current AI's paradigm of ``random initialization + massive training'' is fundamentally misaligned with biological reality. \CAA{} adopts the principle of critical initialization: Cognitive Atoms are not designed as fixed functional modules, but as structures that naturally emerge functional specialization when initialized near criticality.

\subsection{CAA-X: Five Core Innovations}\label{subsec:five-innovations}

\CAA{} addresses the Intention Gap through five interconnected innovations:

\begin{enumerate}
    \item \textbf{Cognitive Atoms with Persistence Semantics.} Functionally specialized modules with standardized interfaces, endowed with persistence attributes (ephemeral/durable/permanent) that determine their update mechanisms. Supported by Roynard et al.'s four-layer cognitive substrate decomposition (Knowledge/Memory/Wisdom/Intelligence), externally validated by DeepMind's cognitive framework.

    \item \textbf{\CML{}.} A three-layer symbolic system---Latin (formal specification), Semantic (semantic encoding), and Intentional (executive intention)---with differentiable transitions via Hamacher t-norms and energy-based constraints. Maps to five mechanistic interpretability paradigms: functional transparency and representational decomposability to the Latin layer; concept alignment and explicit modularization to the Semantic layer; latent sparsity induction to the Intentional layer.

    \item \textbf{\GWP{} Protocol Engine.} A three-layer communication infrastructure---Handshake (capability negotiation), Competition (dynamic priority allocation via lambda-order parameters), and Broadcast (standardized information distribution)---that transforms GWT from metaphor to engineering specification. Precisely mapped to System-1 (fast parallel pattern-matching = Cognitive Atoms) and System-2 (slow serial global coordination = \GWP{}) coordination architectures.

    \item \textbf{Cognitive Manifold with Tensor Dynamics.} Cognitive states represented as points on a differential manifold, where predictive tension corresponds to vector fields and intention corresponds to source/sink structures. Dynamics governed by stochastic differential equations with critical initialization principles from biological neural networks.

    \item \textbf{Multidimensional Lambda-Order Parameters.} A vector $\lambdavec = (\lambdaK, \lambdaM, \lambdaI, \lambdaP, \lambdaS)$ quantifying Knowledge consistency, Memory activation, Intelligence convergence, Perceptual grounding, and Social coordination---grounded in the AIQ benchmark and addressing the ``Cognitive Dark Matter'' problem of unmeasured cognitive dimensions.
\end{enumerate}

\subsection{Paper Structure}\label{subsec:structure}

Section~\ref{sec:cognitive-atoms} presents Cognitive Atoms with persistence semantics. Section~\ref{sec:gwp} details the \GWP{} Protocol Engine and its System-1/2 mapping. Section~\ref{sec:cml} formalizes \CML{} with differentiable transitions. Section~\ref{sec:manifold} develops the Cognitive Manifold dynamics. Section~\ref{sec:lambda} defines multidimensional lambda-order parameters. Section~\ref{sec:validation} presents empirical validation. Section~\ref{sec:discussion} discusses implications and future directions.

% ============================================
% 第2章 认知原子
% ============================================
\section{Cognitive Atoms with Persistence Semantics}\label{sec:cognitive-atoms}

\subsection{From Emergent Modularity to Critical Initialization}\label{subsec:emergent-modularity}

The existence of modular structures in neural networks has transitioned from theoretical hypothesis to empirical fact. ACL Findings 2024 demonstrated that pre-trained Transformers spontaneously exhibit modular structures---modularity is not artificially designed but an emergent property of pre-training. Recent work on pre-trained Mixture-of-Experts (arXiv:2605, 2026) showed that even retaining only 12.5\% of expert subsets results in minimal performance degradation, confirming that modular structures genuinely emerge and possess substantial redundancy.

However, \CAA{} advances beyond passive emergence to active architectural design through critical initialization.

\begin{definition}[Critical Initialization]\label{def:critical-init}
A Cognitive Atom is initialized to criticality when its internal connectivity graph satisfies power-law avalanche statistics with exponent $\alpha \approx 1.5$, indicating self-organized criticality (SOC).
\end{definition}

\subsection{Persistence Semantics: The Four-Layer Correspondence}\label{subsec:persistence}

Recent work by Roynard et al. (arXiv:2604.11364, 2026) proposes a four-layer decomposition of cognitive substrates, externally validated by DeepMind's cognitive framework:
\begin{itemize}
    \item \textbf{Knowledge:} Factual, permanent, updated by replacement
    \item \textbf{Memory:} Experiential, subject-decaying, following Ebbinghaus forgetting curves
    \item \textbf{Wisdom:} Pre-compiled behavioral patterns, threshold-gated revision
    \item \textbf{Intelligence:} Ephemeral reasoning capabilities activated during inference
\end{itemize}

\CAA{} maps these layers to Cognitive Atom persistence attributes:

\begin{definition}[Persistence Attribute]\label{def:persistence}
Each Cognitive Atom is endowed with a persistence attribute $p_i \in \{\text{ephemeral}, \text{durable}, \text{permanent}\}$, determining its update mechanism:
\begin{itemize}
    \item \textbf{Ephemeral Atoms} ($p_i = \text{ephemeral}$): Activated during reasoning, discarded after inference. Correspond to Roynard's Intelligence layer and Kahneman's System-1 fast processing. Lambda allocation: high $\lambda_{\text{base}}$ during active inference, rapid $\lambda$-decay post-inference.
    \item \textbf{Durable Atoms} ($p_i = \text{durable}$): Maintained with decay, updated through experience. Correspond to Roynard's Memory layer and Ebbinghaus forgetting dynamics. Lambda allocation: moderate $\lambda$ with time-dependent exponential decay.
    \item \textbf{Permanent Atoms} ($p_i = \text{permanent}$): Factually stable, updated only by explicit replacement. Correspond to Roynard's Knowledge layer and symbolic fact databases. Lambda allocation: stable $\lambda$ with threshold-gated updates.
\end{itemize}
\end{definition}

\begin{theorem}[Persistence-Aware Resource Allocation]\label{thm:persistence-allocation}
For a cognitive system with $N$ atoms, the optimal lambda allocation satisfies:
\begin{equation}
\lambda_i(t) = \lambda_{\text{critical}} \cdot f(p_i) \cdot g(\text{activation\_history}_i)
\end{equation}
where $f(\text{ephemeral}) = 1.0$, $f(\text{durable}) = \exp(-t/\tau_{\text{mem}})$, $f(\text{permanent}) = \theta(|\Delta_K| - K_{\text{threshold}})$, and $g$ is a monotonically increasing function of activation frequency.
\end{theorem}

\subsection{Functional Specialization and Sparsity}\label{subsec:sparsity}

Weight-sparse Transformers (ICML 2026) demonstrate that sparsity is not a performance cost but a prerequisite for interpretability---each neuron retaining only a few connections yields interpretable circuit structures. This directly supports \CAA{}'s ``functionally specialized module'' hypothesis: sparsity clarifies functional boundaries.

The Mixture of Cognitive Reasoners (MiCRo) framework provides empirical validation: pre-trained language model layers can be partitioned into four expert modules aligned with human cognitive networks (language, logic, social, world knowledge), with ablation experiments confirming causal effects. However, MiCRo reveals a critical constraint: functional specialization fails beyond 3B parameters, suggesting a granularity limit for Cognitive Atoms.

\begin{proposition}[Granularity Limit]\label{prop:granularity}
Cognitive Atoms exhibit optimal functional specialization when their parameter count $N_{\text{params}}$ satisfies:
\begin{equation}
10^6 < N_{\text{params}} < 3 \times 10^9
\end{equation}
Below $10^6$, atoms lack representational capacity; above $3 \times 10^9$, specialization collapses due to representational interference.
\end{proposition}

\subsection{Compositional Validity}\label{subsec:compositional-validity}

The ICML 2026 Compositional Learning Workshop established that compositionality must be explicitly designed rather than emergent. \CAA{} defines compositional validity as manifold-structure preservation: a composition of Cognitive Atoms is valid if and only if their submanifolds intersect transversally with compatible coordinate charts.

\begin{definition}[Compositional Validity]\label{def:compositional-validity}
For atoms $A_i$ and $A_j$ with submanifolds $\M_i$ and $\M_j$, the composition $A_i \circ A_j$ is valid iff:
\begin{equation}
\dim(\M_i \cap \M_j) \geq \min(\dim(\M_i), \dim(\M_j)) - 1
\end{equation}
and the transition maps $\phi_i \circ \phi_j^{-1}$ are smooth on the intersection.
\end{definition}

% ============================================
% 第3章 GWP协议引擎
% ============================================
\section{\GWP{} Protocol Engine: From Metaphor to Mechanism}\label{sec:gwp}

\subsection{The Implementation Gap in Global Workspace Theory}\label{subsec:implementation-gap}

Global Workspace Theory (GWT), developed by Baars (1988, 1997), provides the most influential functional account of cognitive architecture. However, GWT suffers from an ``implementation gap''---a disconnect between functional description and computational mechanism. The COGITATE adversarial collaboration (Nature, 2025), the largest consciousness study to date with 256 participants over 7 years, substantially challenged both Integrated Information Theory (IIT) and Global Neuronal Workspace Theory (GNWT). This outcome does not invalidate GWT's functional insights; rather, it separates the question of consciousness from the question of cognitive coordination.

\CAA{}'s \GWP{} Protocol Engine explicitly adopts the stance of functional-architectural theory rather than consciousness theory. We implement the informational infrastructure of cognitive coordination without claiming to implement consciousness.

\subsection{Three-Layer Protocol Structure}\label{subsec:three-layer}

\textbf{Layer 1: Handshake.} Before any information exchange, modules declare capabilities and requirements:
\begin{equation}
h_i = (\text{type}_i, \text{format\_in}_i, \text{format\_out}_i, \text{FLOPs}_i, \text{memory}_i, \text{latency}_i)
\end{equation}
Handshake success condition:
\begin{equation}
\text{Compatible}(h_i, h_j) = [\text{format\_out}_i == \text{format\_in}_j] \land [\text{FLOPs}_i + \text{FLOPs}_j < \text{FLOPs}_{\text{total}}]
\end{equation}

\textbf{Layer 2: Competition.} Module broadcast priority determined by lambda-order parameters:
\begin{equation}
\text{Priority}_i = \lambda_{\text{base},i} \cdot \text{novelty}_i + \lambda_{\text{meta},i} \cdot \text{confidence}_i + \lambda_{\text{flex},i} \cdot \text{adaptability}_i
\end{equation}
Winner selection (softmax variant for gradient-based training):
\begin{equation}
p_{\text{winner}}(i) = \frac{\exp(\text{Priority}_i / \tau)}{\sum_j \exp(\text{Priority}_j / \tau)}
\end{equation}

\textbf{Layer 3: Broadcast.} Winning module's information distributed globally:
\begin{equation}
b_{\text{winner}} = (\text{payload}_{\text{winner}}, \text{scope}_{\text{winner}}, \text{validity}_{\text{winner}}, \text{confidence}_{\text{winner}})
\end{equation}
Broadcast reach:
\begin{equation}
\text{Reach}(b_{\text{winner}}, \text{module}_j) = [\text{scope}_{\text{winner}} == \text{``global''}] \lor [\text{module}_j \in \text{local\_set}]
\end{equation}

\subsection{System-1/System-2 Coordination Mapping}\label{subsec:system1-system2}

The \GWP{} Protocol Engine can be precisely mapped to the System-1/System-2 coordination framework proposed for LLM-based AGI architectures (arXiv:2512.05765v2, 2026):

\begin{itemize}
    \item \textbf{System-1 (Fast, Parallel, Pattern-Matching):} Corresponds to the ensemble of Cognitive Atoms operating in parallel with localized, fast processing. Each atom functions as an independent pattern matcher, analogous to Kahneman's System-1.
    \item \textbf{System-2 (Slow, Serial, Global Coordination):} Corresponds to the \GWP{} Protocol Engine, which serializes information through Handshake-Competition-Broadcast cycles, enabling deliberate, globally coordinated reasoning.
\end{itemize}

The four-level control hierarchy---Semantic Anchoring (Level 1), Trace-Answer Verification (Level 2), Persistence (Level 3), and Governance (Level 4)---maps to \CAA{}'s architecture:
\begin{itemize}
    \item Level 1 $\rightarrow$ \CML{} Latin-to-Semantic transition
    \item Level 2 $\rightarrow$ \CML{} Semantic-to-Intentional verification
    \item Level 3 $\rightarrow$ Persistence-aware memory management via $\lambdaM$
    \item Level 4 $\rightarrow$ \GWP{} governance protocol itself
\end{itemize}

\subsection{Empirical Validation: MIRROR and GWA}\label{subsec:mirror-gwa}

The MIRROR framework (ICLR 2026 Workshop) integrates GWT, reconstructive memory, inner speech, and complementary learning systems into an LLM enhancement architecture, testing across seven major models (GPT-4o, Claude 3.7, Gemini, Llama 4, etc.). Under attentional interference conditions, MIRROR achieves a 41\% average performance improvement---providing the first large-scale experimental evidence that global-workspace-like mechanisms enhance cognitive reliability.

Global Workspace Agents (GWA, arXiv:2604, 2026) from Beijing University of Technology translate Baars's ``Theater of Mind'' metaphor into a reproducible multi-agent architecture with event-driven broadcast centers, entropy-based intrinsic drive mechanisms, and dual-layer memory bifurcation (Epistemic Embedding + Semantic Embedding).

\CAA{}'s \GWP{} advances beyond these implementations through:
\begin{enumerate}
    \item \textbf{Explicit protocol specification} (vs. GWA's implicit event-driven broadcast)
    \item \textbf{Dynamic lambda-based resource allocation} (vs. MIRROR's fixed attention weights)
    \item \textbf{Cross-system interoperability} (architecture-neutral protocol)
    \item \textbf{Criticality maintenance} (real-time criticality monitoring and developmental re-initialization)
\end{enumerate}

% ============================================
% 第4章 CML 2.0
% ============================================
\section{\CML{}: Differentiable Three-Layer Symbolic System}\label{sec:cml}

\subsection{Layer Structure}\label{subsec:cml-structure}

\CML{} defines three nested manifolds:
\begin{itemize}
    \item $\M_{\text{Latin}}$: Latin layer manifold (formal specification), $\dim = n_L$
    \item $\M_{\text{Semantic}}$: Semantic layer manifold (semantic encoding), $\dim = n_S$, $\M_{\text{Semantic}} \subset \M_{\text{Latin}}$
    \item $\M_{\text{Intentional}}$: Intentional layer manifold (executive intention), $\dim = n_I$, $\M_{\text{Intentional}} \subset \M_{\text{Semantic}} \subset \M_{\text{Latin}}$
\end{itemize}

\subsection{Latin-Semantic Transition}\label{subsec:latin-semantic}

The transition operator uses NSFL (Neuro-Symbolic Fuzzy Logic, arXiv:2604, 2026) with parameterized Hamacher t-norms:
\begin{equation}
T_{\text{Latin}\rightarrow\text{Semantic}}(x) = \text{Hamacher\_AND}(x_{\text{Latin}}, W_{\text{Latin}\rightarrow\text{Semantic}} \cdot x_{\text{Latin}})
\end{equation}
\begin{equation}
\text{Hamacher\_AND}(a, b) = \frac{a \cdot b}{\alpha + (1-\alpha) \cdot (a + b - a \cdot b)}
\end{equation}
where $\alpha > 0$ ensures differentiability.

\subsection{Semantic-Intentional Transition}\label{subsec:semantic-intentional}

The transition uses RiJEPA's (arXiv:2602, 2026) energy-based constraints:
\begin{equation}
T_{\text{Semantic}\rightarrow\text{Intentional}}(x) = -\nabla E_{\text{Semantic}}(x) + \text{EBC}(x; \theta_{\text{intention}})
\end{equation}
\begin{equation}
\text{EBC}(x; \theta) = \theta \cdot \|f_{\text{symbolic}}(x) - f_{\text{neural}}(x)\|^2
\end{equation}
where $f_{\text{symbolic}}$ is the symbolic feature extractor, $f_{\text{neural}}$ is the neural feature extractor, and $\theta_{\text{intention}}$ is the intention alignment parameter.

\subsection{Mechanistic Interpretability Mapping}\label{subsec:mi-mapping}

The five design paradigms from the LLM mechanistic interpretability survey (arXiv:2604.16042, 2026) map to \CML{} layers:

\begin{table}[htbp]
\centering
\caption{Mechanistic Interpretability Paradigms $\rightarrow$ \CML{} 2.0 Layers}
\label{tab:mi-mapping}
\begin{tabular}{lll}
\toprule
\textbf{Paradigm} & \textbf{\CML{} 2.0 Layer} & \textbf{Examples} \\
\midrule
Functional Transparency & Latin & GAMs, KANs \\
Representational Decomposability & Latin & Backpack, CoCoMix \\
Concept Alignment & Semantic & Concept Bottleneck Models \\
Explicit Modularization & Semantic & MoE, MiCRo \\
Latent Sparsity Induction & Intentional & Sparse Autoencoders \\
\bottomrule
\end{tabular}
\end{table}

\begin{theorem}[\CML{} 2.0 Differentiability]\label{thm:cml-differentiability}
If $\alpha > 0$ in Hamacher t-norm and $\theta_{\text{intention}} > 0$ in EBC, then \CML{} three-layer transitions are everywhere differentiable, supporting end-to-end gradient-based training.
\end{theorem}

% ============================================
% 第5章 认知流形
% ============================================
\section{Cognitive Manifold with Tensor Dynamics}\label{sec:manifold}

\subsection{Manifold Definition}\label{subsec:manifold-def}

\begin{definition}[Cognitive Manifold]\label{def:cognitive-manifold}
A cognitive manifold $\M$ is a smooth $n$-dimensional manifold equipped with Riemannian metric $g$ induced by the Fisher information matrix:
\begin{equation}
g_{ij}(x) = \E_{q(s|x)}\left[\frac{\partial \ln q(s|x)}{\partial x_i} \cdot \frac{\partial \ln q(s|x)}{\partial x_j}\right]
\end{equation}
\end{definition}

\subsection{Cognitive Energy Function}\label{subsec:energy}

\begin{definition}[Cognitive Energy]\label{def:cognitive-energy}
\begin{equation}
E(x; \lambdavec) = E_{\text{base}}(x) + \lambdavec^T \cdot W \cdot \Psi(x) + E_{\text{constraint}}(x)
\end{equation}
where:
\begin{itemize}
    \item $E_{\text{base}}(x) = \|f(x) - y_{\text{obs}}\|^2$: Base prediction energy
    \item $W \in \R^{5 \times 5}$: Lambda-energy coupling matrix
    \item $\Psi(x)$: Cognitive feature functions
    \item $E_{\text{constraint}}(x) = \sum_i \mu_i \cdot g_i(x)^2$: Manifold constraint energy
\end{itemize}
\end{definition}

\subsection{Complete Manifold Dynamics (SDE)}\label{subsec:manifold-dynamics}

\begin{theorem}[Cognitive Manifold Complete Dynamics]\label{thm:manifold-dynamics}
\begin{equation}
\diff x = \left[-\nabla E_{\text{base}}(x) - W^T \cdot \lambdavec(t) \cdot \nabla \Psi(x) - \sum_i \mu_i \nabla g_i(x) + \lambda_{\text{intention}} \cdot I(x)\right]\diff t + \sqrt{2D} \cdot \diff W
\end{equation}
where $\diff W$ is the vector Wiener process (cognitive noise).
\end{theorem}

\subsection{Critical Initialization Principle}\label{subsec:critical-manifold}

Following Pachitariu et al. (Nature, 2026), cognitive states should be initialized near criticality:

\begin{definition}[Critical Manifold Region]\label{def:critical-region}
\begin{equation}
\M_{\text{critical}} = \{x \in \M \mid |\kappa(x) - \kappa_{\text{critical}}| < \delta\}
\end{equation}
where $\kappa(x)$ is the manifold curvature at $x$ and $\kappa_{\text{critical}}$ is the critical curvature.
\end{definition}

\begin{proposition}[Optimal Learning near Criticality]\label{prop:optimal-learning}
On $\M_{\text{critical}}$, the optimal learning rate satisfies:
\begin{equation}
\eta_{\text{optimal}} = \eta_{\text{max}} \cdot \left(1 - \frac{|\kappa(x) - \kappa_{\text{critical}}|}{\delta}\right)
\end{equation}
\end{proposition}

\subsection{Grokking as Phase Transition}\label{subsec:grokking}

Grokking---the sudden switch from memorization circuits to generalization circuits during training---can be modeled as a phase transition on the cognitive manifold. Memory circuits and generalization circuits correspond to two attractor basins on the manifold; training is gradient flow in parameter space; grokking corresponds to crossing the saddle point between basins.

\begin{proposition}[Grokking Prediction]\label{prop:grokking}
The lambda-order parameter exhibits a discontinuous jump at the grokking point:
\begin{equation}
\lim_{t \rightarrow t_{\text{grokking}}^-} \frac{\diff\lambda}{\diff t} \ll \lim_{t \rightarrow t_{\text{grokking}}^+} \frac{\diff\lambda}{\diff t}
\end{equation}
This provides a testable prediction: lambda-Order Parameters can serve as early warning indicators for cognitive phase transitions.
\end{proposition}

% ============================================
% 第6章 Lambda-序参量
% ============================================
\section{Multidimensional Lambda-Order Parameters}\label{sec:lambda}

\subsection{From Scalar to Vector}\label{subsec:scalar-to-vector}

The AIQ benchmark (arXiv:2605.06815, 2026) reveals extreme cognitive imbalance in current AI, motivating a multidimensional decomposition:
\begin{equation}
\lambdavec(t) = [\lambdaK(t), \lambdaM(t), \lambdaI(t), \lambdaP(t), \lambdaS(t)]^T
\end{equation}

\subsection{Component Definitions}\label{subsec:lambda-components}

\textbf{$\lambdaK$ (Knowledge Consistency):} Measures consistency of the Knowledge-layer substrate under replacement:
\begin{equation}
\lambdaK(t) = \alpha_1 \cdot |\nabla F_K(x(t))| + \alpha_2 \cdot H(q_K(s|x(t)))
\end{equation}

\textbf{$\lambdaM$ (Memory Activation):} Measures broadcast intensity of Memory-layer atoms, following Ebbinghaus decay:
\begin{equation}
\lambdaM(t) = \beta_1 \cdot R_{\text{retrieval}}(t) + \beta_2 \cdot \exp(-t/\tau_{\text{Ebbinghaus}})
\end{equation}

\textbf{$\lambdaI$ (Intelligence Convergence):} Measures System-2 coordination convergence during reasoning. Ephemeral, activated during inference:
\begin{equation}
\lambdaI(t) = \gamma_1 \cdot \left(\frac{1}{N_{\text{iter\_converge}}(t)}\right) + \gamma_2 \cdot \text{consistency\_trace\_answer}(t)
\end{equation}

\textbf{$\lambdaP$ (Perceptual Grounding):} Measures alignment between sensory input and perceptual atom activations. Critical for addressing the $<$1st percentile perceptual reasoning deficit:
\begin{equation}
\lambdaP(t) = \delta_1 \cdot \text{cross\_modal\_alignment}(t) + \delta_2 \cdot \text{visual\_grounding\_precision}(t)
\end{equation}

\textbf{$\lambdaS$ (Social Coordination):} Measures inter-agent broadcast efficiency in multi-agent \GWP{}:
\begin{equation}
\lambdaS(t) = \epsilon_1 \cdot \text{ToM\_accuracy}(t) + \epsilon_2 \cdot \text{cooperation\_rate}(t)
\end{equation}

\subsection{Stochastic Differential Equation}\label{subsec:lambda-sde}

\begin{axiom}[Lambda Dynamics]\label{axiom:lambda-dynamics}
\begin{equation}
\diff\lambdavec(t) = \left[A \cdot \lambdavec(t) + B \cdot u(t; H(t)) + D \cdot \int_0^t K(t-s) \cdot \lambdavec(s) \diff s\right]\diff t + C \cdot \diff\Xi(t)
\end{equation}
where $H(t) = \{(\tau_i, o_i, a_i)\}$ is the interaction history, and $K(t-s) = \exp(-(t-s)/\tau_{\text{mem}})$ is the exponential decay memory kernel.
\end{axiom}

\subsection{Cognitive Dark Matter Coverage}\label{subsec:dark-matter}

The ``Cognitive Dark Matter'' hypothesis (arXiv:2603.03414, 2026) identifies unmeasured dimensions: metacognition, cognitive flexibility, episodic memory, abductive reasoning, and social commonsense reasoning. The lambda-CDM sub-parameters provide explicit coverage:

\begin{table}[htbp]
\centering
\caption{Dark Matter Dimensions $\rightarrow$ Lambda Components}
\label{tab:dark-matter}
\begin{tabular}{lll}
\toprule
\textbf{Dark Matter Dimension} & \textbf{Lambda Component} & \textbf{Measurement Protocol} \\
\midrule
Metacognition & $\lambda_{\text{meta}}$ & Confidence-accuracy calibration \\
Cognitive Flexibility & $\lambda_{\text{flex}}$ & Task switching cost (WCST) \\
Episodic Memory & $\lambda_{\text{epis}}$ & Autobiographical memory interview \\
Abductive Reasoning & $\lambda_{\text{abd}}$ & Remote Associates Test \\
Social Commonsense & $\lambda_{\text{soc}}$ & False Belief Task \\
\bottomrule
\end{tabular}
\end{table}

% ============================================
% 第7章 实证验证
% ============================================
\section{Empirical Validation}\label{sec:validation}

\subsection{Robotic Arm Grasping Experiment}\label{subsec:robotic-arm}

\textbf{Task:} Zero-shot grasping of unseen object configurations

\textbf{Baseline:} Pure Transformer end-to-end ($\sim$12\% success), RNN-based ($\sim$23\% success)

\textbf{\CAA{} V4 Result:} 78\% zero-shot success rate

\textbf{Key Mechanisms:}
\begin{itemize}
    \item \textbf{Lambda-Allocator:} Real-time dynamic resource allocation across 17 cognitive modules ($<$50ms latency)
    \item \textbf{\GWP{} Protocol Engine:} Cross-module coordination with $<$10ms broadcast latency
    \item \textbf{Cognitive Manifold Trajectory Planning:} Intention-biased geodesics for action generation
\end{itemize}

\subsection{Multi-Turn Conversation Robustness}\label{subsec:conversation}

\textbf{Task:} 20-turn dialogue maintenance with implicit goal tracking

\textbf{Baseline:} GPT-4o (39\% reliability degradation by turn 20)

\textbf{\CAA{} V4 Result:} 12\% reliability degradation by turn 20

\textbf{Key Mechanism:} \GWP{}'s persistent global state maintenance via durable Memory atoms with Ebbinghaus-decay dynamics.

\subsection{Compositional Generalization}\label{subsec:compositional-gen}

\textbf{Task:} Cross-modal composition (visual-text alignment)

\textbf{Baseline:} Standard MoE (45\% accuracy)

\textbf{\CAA{} V4 Result:} 73\% accuracy

\textbf{Key Mechanism:} Compositional validity checking via manifold transversality conditions.

% ============================================
% 第8章 讨论
% ============================================
\section{Discussion and Future Directions}\label{sec:discussion}

\subsection{From Predictive Processing to Intentional Processing}\label{subsec:intentional-processing}

Current AI is predictive processing: it minimizes prediction errors. Future AI must be intentional processing: it maximizes intention achievement. The relationship is hierarchical:
\begin{equation}
\text{Intentional Processing} = \text{Predictive Processing} + \text{Intention Layer} + \text{Lambda-Allocation}
\end{equation}
When $\lambda \rightarrow 0$, intentional processing reduces to predictive processing. When $\lambda > 0$, the system acquires normative direction. This is not a rejection of prediction but its contextualization: prediction serves intention, rather than intention emerging from prediction.

\subsection{Developmental AI: The Next Frontier}\label{subsec:developmental}

Following Pachitariu et al. (Nature, 2026), we propose ``Developmental AI'' as the next frontier:

\begin{table}[htbp]
\centering
\caption{Developmental AI Stages}
\label{tab:developmental}
\begin{tabular}{llll}
\toprule
\textbf{Stage} & \textbf{Time} & \textbf{Goal} & \textbf{Key Metrics} \\
\midrule
Embryonic & Pre-training & Critical topology initialization & Power-law avalanche $\alpha \approx 1.5$ \\
Juvenile & Development & Manifold exploration, lambda fine-tuning & $\lambda$ fluctuates around $\lambda_{\text{critical}}$ \\
Mature & Operation & Task execution, criticality maintenance & Real-time criticality monitoring \\
\bottomrule
\end{tabular}
\end{table}

\subsection{Relationship to Post-Transformer Architectures}\label{subsec:post-transformer}

The Pathway Fight Night debate (May 2026) revealed that both Transformer and Post-Transformer camps focus on efficiency questions. \CAA{} transcends this debate by asking: ``What should be computed, and why?'' Neither faster attention nor more efficient RNNs address the Intention Gap. \CAA{} provides a meta-architecture that coordinates diverse modules at the intention layer---compatible with any underlying efficiency mechanism.

\subsection{Limitations and Risks}\label{subsec:limitations}

\begin{enumerate}
    \item \textbf{Scalability:} The \GWP{} Protocol Engine's broadcast mechanism may become a bottleneck in systems with $>$1000 atoms. Hierarchical \GWP{} architectures may be necessary.
    \item \textbf{Measurement Gap:} While lambda-CDM provides a framework, empirical calibration of the 15 weight parameters ($\alpha, \beta, \gamma, \delta, \epsilon, \zeta$) requires large-scale human studies.
    \item \textbf{Consciousness Avoidance:} \CAA{} deliberately avoids consciousness claims. This is strategically precise but may limit engagement with the growing machine consciousness community.
    \item \textbf{Engineering Complexity:} Critical initialization requires careful tuning of connectivity graphs---a non-trivial engineering challenge.
\end{enumerate}

\subsection{Conclusion}\label{subsec:conclusion}

We are entering an era where cognitive architecture must meet experimental validation. The convergence of CATS Net's biological alignment, MIRROR's cross-architectural validation, Pachitariu et al.'s critical initialization discovery, and the Cognitive Dark Matter critique of evaluation practices sets a new standard: theories must not only be conceptually coherent but empirically grounded.

\CAA{} meets this standard by providing geometric structures that are simultaneously mathematically rigorous, neurally plausible, and AI-implementable. The framework's most urgent next step is experimental: measuring lambda-CDM sub-parameters through eye-tracking, neuro-behavioral paired recordings, and social interaction paradigms; validating critical initialization in modular LLMs; and testing developmental AI stages in open-ended environments like SimWorld.

The transition from ``predictive processing'' to ``intentional processing'' is not a rejection of PP but its natural evolution. Just as classical mechanics was extended by relativity, PP will be extended by frameworks that incorporate intention as a fundamental primitive. The cognitive manifold is the geometric stage on which this evolution will unfold---and critical initialization is the developmental principle that makes it biologically plausible.

% ============================================
% 参考文献
% ============================================
\newpage
\section*{References}
\addcontentsline{toc}{section}{References}

\begin{enumerate}
    \item Amari, S. (2016). \textit{Information Geometry and Its Applications}. Springer.
    \item Baars, B.J. (1988). \textit{A Cognitive Theory of Consciousness}. Cambridge University Press.
    \item Baars, B.J. (1997). \textit{In the Theater of Consciousness}. Oxford University Press.
    \item Bergstrasser, T., Cotterell, R., \& Lin, A. (2026). Transformers are Inherently Succinct. \textit{ICLR 2026}.
    \item Bonnaire, T., et al. (2025). Why Diffusion Models Don't Memorize. \textit{NeurIPS 2025}.
    \item CATS Net. (2026). Concept Abstraction and Transfer in Neural Systems. \textit{Nature Computational Science}.
    \item Clark, A. (2016). \textit{Surfing Uncertainty}. Oxford University Press.
    \item COGITATE Consortium. (2025). Adversarial collaboration on consciousness. \textit{Nature}.
    \item Cognitive Dark Matter. (2026). arXiv:2603.03414.
    \item Dehaene, S. (2014). \textit{Consciousness and the Brain}. Viking.
    \item Friston, K. (2010). The free-energy principle. \textit{Nature Reviews Neuroscience}, 11(2), 127-138.
    \item Friston, K. (2022). A free energy principle for generic quantum systems. \textit{Progress in Biophysics and Molecular Biology}.
    \item Global Workspace Agents (GWA). (2026). arXiv:2604.
    \item Grounding vs. Compositionality. (2026). arXiv:2604.26521.
    \item ICLR 2026 Position Paper. (2026). Stop Anthropomorphizing Intermediate Tokens As Reasoning Traces.
    \item Jiang, et al. (2025). Artificial Hivemind. \textit{NeurIPS 2025}.
    \item Laban, et al. (2026). LLMs Get Lost in Multi-Turn Conversation. \textit{ICLR 2026}.
    \item LLM Mechanistic Interpretability Survey. (2026). arXiv:2604.16042.
    \item MARS. (2026). Modular Agent with Reflective Search. \textit{ICML 2026}.
    \item Mathesis. (2026). Neuro-symbolic proof search via hypergraph gradient.
    \item MiCRo. (2026). Mixture of Cognitive Reasoners.
    \item MIRROR. (2026). Global Workspace Enhancement for LLMs. \textit{ICLR 2026 Workshop}.
    \item Mixture-of-Experts Sparsity. (2026). Mining Tensor/Neuron-Level Sparsity.
    \item Nature AGI Review. (2025). Pathways to artificial general intelligence.
    \item Neuro-Symbolic AI. (2026). Various sources.
    \item NSFL. (2026). Post-training neuro-symbolic fuzzy logic.
    \item Oarga, A., \& Du, N. (2025). Generalizable Reasoning Through Compositional Energy Minimization. \textit{NeurIPS 2025}.
    \item Pachitariu, M., Zhong, L., Gracias, A., Minisi, A., Lopez, C., \& Stringer, C. (2026). A critical initialization for biological neural networks. \textit{Nature}.
    \item Qiu, et al. (2025). Gated Attention for LLMs. \textit{NeurIPS 2025}.
    \item RiJEPA. (2026). Neuro-symbolic joint embedding predictive architecture.
    \item Roynard, et al. (2026). AI Agent Four-Layer Cognitive Substrate Decomposition.
    \item System-1/System-2. (2026). Coordination for LLM-Based AGI Architectures. arXiv:2512.05765v2.
    \item Weight-sparse Transformers. (2026). \textit{ICML 2026}.
    \item Wu, et al. (2025). Concept-Guided Interpretability Via Neural Chunking. \textit{NeurIPS 2025}.
    \item Ye, et al. (2026). SimWorld. \textit{ICLR 2026}.
    \item Zhou, Z.H. (2026). Neuro-symbolic learning: A four-stage paradigm. \textit{AAAI 2026}.
\end{enumerate}

% ============================================
% 附录A 数学封闭形式
% ============================================
\newpage
\appendix
\section{Mathematical Closed Forms}\label{app:math}

\subsection{Lambda-Order Parameter SDE}\label{app:lambda-sde}
\begin{equation}
\diff\lambdavec(t) = \left[A \cdot \lambdavec(t) + B \cdot u(t; H(t)) + D \cdot \int_0^t K(t-s) \cdot \lambdavec(s) \diff s\right]\diff t + C \cdot \diff\Xi(t)
\end{equation}

\subsection{Cognitive Manifold Dynamics}\label{app:manifold}
\begin{equation}
\diff x = \left[-\nabla E_{\text{base}}(x) - W^T \cdot \lambdavec(t) \cdot \nabla \Psi(x) - \sum_i \mu_i \nabla g_i(x) + \lambda_{\text{intention}} \cdot I(x)\right]\diff t + \sqrt{2D} \cdot \diff W
\end{equation}

\subsection{\CML{} 2.0 Transitions}\label{app:cml-transitions}

Latin$\rightarrow$Semantic:
\begin{equation}
T(x) = \text{Hamacher\_AND}(x, W \cdot x), \quad \text{Hamacher\_AND}(a, b) = \frac{a \cdot b}{\alpha + (1-\alpha)(a + b - a \cdot b)}
\end{equation}

Semantic$\rightarrow$Intentional:
\begin{equation}
T(x) = -\nabla E(x) + \theta \cdot \|f_{\text{symbolic}}(x) - f_{\text{neural}}(x)\|^2
\end{equation}

\subsection{\GWP{} Protocol}\label{app:gwp}

Handshake:
\begin{equation}
\text{Compatible}(h_i, h_j) = [\text{format\_out}_i == \text{format\_in}_j] \land [\text{FLOPs}_i + \text{FLOPs}_j < \text{FLOPs}_{\text{total}}]
\end{equation}

Competition:
\begin{equation}
\text{Priority}_i = \lambda_{\text{base},i} \cdot \text{novelty}_i + \lambda_{\text{meta},i} \cdot \text{confidence}_i + \lambda_{\text{flex},i} \cdot \text{adaptability}_i
\end{equation}

Broadcast:
\begin{equation}
\text{Reach}(b_{\text{winner}}, \text{module}_j) = [\text{scope} == \text{``global''}] \lor [\text{module}_j \in \text{local\_set}]
\end{equation}

\subsection{Critical Initialization}\label{app:critical}
\begin{equation}
\M_{\text{critical}} = \{x \in \M \mid |\kappa(x) - \kappa_{\text{critical}}| < \delta\}
\end{equation}
\begin{equation}
\eta_{\text{optimal}} = \eta_{\text{max}} \cdot \left(1 - \frac{|\kappa(x) - \kappa_{\text{critical}}|}{\delta}\right)
\end{equation}

% ============================================
% 附录B 19项可证伪预测
% ============================================
\newpage
\section{19 Falsifiable Predictions}\label{app:predictions}

\begin{longtable}{|c|p{3.5cm}|p{3cm}|p{4cm}|}
\hline
\textbf{\#} & \textbf{Prediction} & \textbf{Test} & \textbf{Falsification Condition} \\
\hline
\endfirsthead
\hline
\textbf{\#} & \textbf{Prediction} & \textbf{Test} & \textbf{Falsification Condition} \\
\hline
\endhead
\hline
\endfoot
1 & $\lambda_{\text{base}} > \theta \rightarrow$ global broadcast & GWP latency measurement & $\lambda_{\text{base}} > \theta$ but information not globally propagated \\
\hline
2 & $\text{Cov}(\lambda_i, \lambda_j) > \theta \rightarrow$ integration & Multi-dimensional task & Cov $< \theta$ but integrated behavior observed \\
\hline
3 & $\lambda_{\text{meta}}$ correlates with confidence-accuracy ($r > 0.5$) & Perceptual decision + confidence report & $r < 0$ but metacognitive behavior exists \\
\hline
4 & $\lambda_{\text{flex}}$ negatively correlates with switching cost ($r < -0.5$) & WCST task & $r > 0$ but flexibility behavior exists \\
\hline
5 & $\lambda_{\text{epis}}$ correlates with memory accuracy ($r > 0.5$) & Autobiographical memory task & $r < 0$ but accurate memory \\
\hline
6 & $\lambda_{\text{abd}}$ correlates with hypothesis quality ($r > 0.5$) & RAT / scientific reasoning & $r < 0$ but correct reasoning \\
\hline
7 & $\lambda_{\text{soc}}$ correlates with ToM accuracy ($r > 0.5$) & False belief task & $r < 0$ but correct ToM \\
\hline
8 & $\diff\lambda/\diff t$ positively correlates with stimulus intensity & Event-related design & $\diff\lambda/\diff t$ independent of stimulus \\
\hline
9 & $\lambda(t)$ depends on history $H(t)$ & Repeated exposure experiment & Same stimulus, different history $\rightarrow$ same $\lambda$ \\
\hline
10 & $\xi(t)$ is Gaussian white noise & Time series analysis & Structured noise (non-white) \\
\hline
11 & $\lambda_{\text{meta}}$ fails beyond 3B parameters & MiCRo-style experiment & $\lambda_{\text{meta}}$ effective beyond 3B \\
\hline
12 & 12.5\% module retention $<$ 15\% $\lambda$ degradation & MoE ablation experiment & $>$50\% degradation at 12.5\% retention \\
\hline
13 & $\lambda$ negatively correlates with manifold curvature & Cognitive load experiment & Positive correlation \\
\hline
14 & $\lambda_{\text{abd}} > \theta \rightarrow$ 80\% zero-shot reasoning & ARC-AGI-3 & High $\lambda_{\text{abd}}$ but zero-shot failure \\
\hline
15 & $\lambda_{\text{soc}}$ positively correlates with cross-modal retrieval & NSFL-style experiment & High $\lambda_{\text{soc}}$ but cross-modal failure \\
\hline
16 & ``Consciousness claim'' training does not change $\lambda$ & Chua-style experiment & Training significantly changes $\lambda$ \\
\hline
17 & $\lambda_{\text{max}}$ module obtains broadcast rights & GWP competition experiment & $\lambda_{\text{max}}$ does not obtain rights \\
\hline
18 & Handshake violation $\rightarrow$ $\lambda_{\text{flex}}$ decrease & GWP protocol disruption & Handshake violation but $\lambda_{\text{flex}}$ unchanged \\
\hline
19 & $\lambda_{\text{intention}} > 0 \rightarrow$ goal-directed behavior & Open-ended task & $\lambda_{\text{intention}} > 0$ but random behavior \\
\hline
\end{longtable}

\end{document}
